\documentclass[a4paper,twoside,openright,11pt]{book}
%\documentclass[a4paper,twoside,openrigh,headinclude,footexclude,mpexclude,DIV11,12pt,BCOR15mm,listsleft,bibtotoc]{book}



% Pakete
\usepackage[T1]{fontenc}

\usepackage[utf8]{inputenc} % Editor auf utf8 stellen
% z.B. bei texlipse in Eclipse unter Windows->Preferences,
% dann General->Workspace unten bei ``Text file e	ncoding''.

\usepackage[german]{babel}

\usepackage{german} % unter anderem für bessere Silbentrennung

\usepackage{listings}
\usepackage{amssymb}
\usepackage{amsmath}
\usepackage{graphicx}
\usepackage{url}
\usepackage{hyperref}
\usepackage{xspace}

%avoid overfull hbox problem
\setlength{\emergencystretch}{3em}

%avoid underfull hbox problem in bibliography
\usepackage{etoolbox}
% Variant A
\apptocmd{\sloppy}{\hbadness 2000\relax}{}{}
% Variant B
% \apptocmd{\thebibliography}{\raggedright}{}{}

\hyphenpenalty=5000
\tolerance=1000



% standard scaling factor for graphics, to support homogenous looks.
\newcommand{\stdScaleFactor}{0.45}

% SE Titelseite und Erklaerung
\usepackage{sethesis}


% Angaben bez\UTF{00FC}glich Universitaet und Fachbereich
\uniname{Gottfried Wilhelm\break Leibniz Universität Hannover}
\fachgebiet{Fachgebiet Software Engineering}
\institut{Institut für Praktische Informatik}
\fachbereich{Fakultät für Elektrotechnik und Informatik}
\studiengang{Informatik}
% Angaben zur Arbeit
\author{<Vorname Nachname>}
\title{<Titel der Arbeit>}
\englishtitle{<Englischer Titel der Arbeit>}
\ort{Hannover}
\thesistype{<Bachelorarbeit>}
%\thesistype{Masterarbeit}
% Pruefer
\erstpru{<Name des Erstprüfers>}
\zweitpru{<Name des Zweitprüfers>}
\betreu{<Name des Betreuers>}
\datum{<Datum der Abgabe>}

% Sonstiges
%\allowdisplaybreaks
%\setlength{\parindent}{1em}
%\setlength{\parskip}{0em}
%\makeatletter
%\renewcommand{\fps@figure}{ht!}
%\renewcommand{\fps@table}{ht!}
%\setlength{\@fptop}{0pt}
%\makeatother




\begin{document}

\frontmatter

\maketitle

\clearpage
\input{affirmation.tex}

\clearpage
\chapter*{Abstract}
\addcontentsline{toc}{chapter}{Abstract}

This document provides comprehensive technical documentation for the L3S-Offshore-2 Frontend, a web-based application for Petri net visualization, manipulation, and simulation. The frontend was developed as part of the L3S-Offshore-2 research project at the L3S Research Center, Leibniz Universität Hannover.

The application is built using Angular and provides the following core capabilities:
\begin{itemize}
    \item \textbf{Visualization:} Interactive rendering of Petri nets from PNML files with automatic layout algorithms (Sugiyama)
    \item \textbf{Manipulation:} Drag-and-drop editing of places, transitions, and arcs with real-time token modification
    \item \textbf{Simulation:} Step-by-step and automated playback of firing sequences with transition highlighting
    \item \textbf{Analysis:} Display of place invariants and transition firing statistics from multi-run simulations
    \item \textbf{Offshore Scheduling:} Dedicated parameter interface for offshore wind farm installation scheduling models
\end{itemize}

The documentation covers the system architecture, key implementation details, deployment procedures, and provides guidance for future development and maintenance.

\clearpage



\tableofcontents

\mainmatter

\chapter{Introduction}
\label{chap:introduction}

% ==============================================================================
% NOTE: This chapter provides an overview of the project context, problem 
% statement, and contributions. Fill in the placeholders with project-specific
% content during the documentation phase.
% ==============================================================================

\section{Project Context}

The L3S-Offshore-2 project is a research initiative at the L3S Research Center focusing on the application of Petri nets for modeling and simulating offshore wind farm installation logistics. Efficient scheduling of vessel operations, component installations, and resource management is critical for minimizing costs and maximizing the utilization of weather windows.

This documentation describes the development of a web-based frontend application that enables researchers and practitioners to:
\begin{itemize}
    \item Visualize complex Petri net models in an interactive browser environment
    \item Modify and extend existing models through a graphical user interface
    \item Execute and analyze simulations to evaluate scheduling strategies
    \item Configure domain-specific parameters for offshore logistics scenarios
\end{itemize}

The frontend integrates with a Python-based backend (Flask) that provides simulation capabilities using the \texttt{gspn4py} library for Generalized Stochastic Petri Nets.

\section{Problem Statement}
\label{sec:problem-statement}

Prior to this project, the research group relied on disconnected tools for Petri net modeling:
\begin{enumerate}
    \item \textbf{Visualization:} Desktop applications (e.g., WoPeD) that could not be easily shared
    \item \textbf{Simulation:} Python scripts executed via command line
    \item \textbf{Parameter Configuration:} Manual editing of configuration files
\end{enumerate}

This fragmentation created barriers for:
\begin{itemize}
    \item Collaboration between researchers
    \item Rapid prototyping and iteration on models
    \item Demonstration of results to external stakeholders
\end{itemize}

\section{Solution Approach}
\label{sec:solution-approach}

The solution is a single-page web application (SPA) that consolidates all workflows into an integrated environment. The key design decisions include:

\begin{itemize}
    \item \textbf{Angular Framework:} Chosen for its component-based architecture and TypeScript support, enabling maintainable and type-safe code
    \item \textbf{Fork of Existing Project:} Instead of building from scratch, the project forked the open-source ``PNML-Frontend'' by i-love-petrinets, which provided a solid foundation for PNML parsing and basic visualization
    \item \textbf{Tab-Based Navigation:} Clear separation of concerns (Build, Code, Simulation, Analyze, Save, Offshore) for different user workflows
    \item \textbf{Containerized Deployment:} Docker-based deployment for easy installation on research servers
\end{itemize}

\section{Contributions}
\label{sec:contributions}

The main contributions of this project are:

\begin{enumerate}
    \item \textbf{Extended Visualization:}
    \begin{itemize}
        \item Viewport-centered zooming and panning with ResizeObserver integration
        \item Automatic layout via Sugiyama algorithm for PNML files without coordinates
        \item Tab-aware UI state management
    \end{itemize}
    
    \item \textbf{Simulation Capabilities:}
    \begin{itemize}
        \item Animated playback of firing sequences from backend simulation
        \item Manual ``Token Game'' mode for step-by-step execution
        \item Multi-run simulation with transition firing frequency statistics
    \end{itemize}
    
    \item \textbf{Offshore-Specific Integration:}
    \begin{itemize}
        \item Dedicated parameter input panel for offshore scheduling models
        \item Integration with backend scheduling endpoints
    \end{itemize}
    
    \item \textbf{Production-Ready Deployment:}
    \begin{itemize}
        \item Docker containerization with nginx reverse proxy
        \item Deployment on L3S research infrastructure
    \end{itemize}
\end{enumerate}

\section{Document Structure}
\label{sec:document-structure}

This documentation is organized as follows:

\begin{description}
    \item[Chapter~\ref{chap:fundamentals}:] Provides background on Petri nets, the PNML format, and relevant web technologies (Angular, TypeScript).
    
    \item[Chapter~\ref{chap:ui-service}:] Introduces the UI Service for managing application state such as the active tab and selected tools.
    
    \item[Chapter~\ref{chap:petri-net-elements}:] Describes the core data model: Place, Transition, Arc, and their TypeScript implementations.
    
    \item[Chapter~\ref{chap:data-service}:] Explains the central data store that manages all Petri net elements and notifies subscribers of changes.
    
    \item[Chapter~\ref{chap:architecture}:] Describes the overall system architecture, component hierarchy, and service dependencies.
    
    \item[Chapter~\ref{chap:implementation}:] Details key implementation aspects including the zoom system, simulation logic, and token reset.
    
    \item[Chapter~\ref{chap:data-io}:] Covers PNML and JSON file import/export functionality.
    
    \item[Chapter~\ref{chap:layout-algorithms}:] Explains the Spring Embedder and Sugiyama layout algorithms.
    
    \item[Chapter~\ref{chap:planning-service}:] Describes backend API integration for simulation and planning.
    
    \item[Chapter~\ref{chap:deployment}:] Covers the Docker-based deployment process and production configuration.
    
    \item[Chapter~\ref{chap:related-work}:] Discusses related tools and how this project compares.
    
    \item[Chapter~\ref{chap:conclusion}:] Summarizes achievements and outlines future work.
    
    \item[Appendix~\ref{chap:appendix-a}:] Contains supplementary material such as API documentation and configuration files.
\end{description}



\chapter{Fundamentals}
\label{chap:fundamentals}

This chapter provides the theoretical and technical background necessary to understand the implementation details described in later chapters. It covers the fundamentals of Petri nets, the PNML exchange format, and the web technologies used in the frontend development.

\section{Petri Nets}
\label{sec:petri-nets}

Petri nets are a mathematical modeling language for the description of distributed systems. Originally introduced by Carl Adam Petri in his 1962 dissertation, they provide a graphical and formal representation of concurrent processes.

\subsection{Basic Definitions}

A Petri net is a directed bipartite graph consisting of:

\begin{itemize}
    \item \textbf{Places} ($P$): Represented as circles, places hold tokens and model conditions or states
    \item \textbf{Transitions} ($T$): Represented as rectangles or bars, transitions model events or actions
    \item \textbf{Arcs} ($F$): Directed edges connecting places to transitions (input arcs) or transitions to places (output arcs)
    \item \textbf{Marking} ($M$): The distribution of tokens across places, representing the system state
\end{itemize}

Formally, a Petri net is defined as a tuple $N = (P, T, F, W, M_0)$ where:
\begin{itemize}
    \item $P$ is a finite set of places
    \item $T$ is a finite set of transitions, $P \cap T = \emptyset$
    \item $F \subseteq (P \times T) \cup (T \times P)$ is the flow relation
    \item $W: F \rightarrow \mathbb{N}^+$ is the arc weight function
    \item $M_0: P \rightarrow \mathbb{N}$ is the initial marking
\end{itemize}

\subsection{Firing Rules}

A transition $t \in T$ is \textbf{enabled} at marking $M$ if and only if:
\[
\forall p \in \bullet t: M(p) \geq W(p, t)
\]

where $\bullet t$ denotes the preset of $t$ (input places). When an enabled transition fires, the marking changes according to:
\[
M'(p) = M(p) - W(p, t) + W(t, p)
\]

\subsection{Extensions: Stochastic Petri Nets}

For modeling real-world systems with timing, Generalized Stochastic Petri Nets (GSPNs) extend basic Petri nets with:
\begin{itemize}
    \item \textbf{Timed transitions}: Fire after a random delay drawn from an exponential distribution
    \item \textbf{Immediate transitions}: Fire instantaneously when enabled (priority over timed transitions)
\end{itemize}

The L3S-Offshore-2 project uses GSPNs to model the stochastic nature of offshore installation operations, including weather-dependent delays.

\section{PNML: Petri Net Markup Language}
\label{sec:pnml}

PNML (Petri Net Markup Language) is an XML-based interchange format for Petri nets, standardized as ISO/IEC 15909-2. It provides a common syntax for storing and exchanging Petri net models across different tools.

\subsection{Document Structure}

A PNML document has the following hierarchical structure:
\begin{lstlisting}[language=XML, caption={Basic PNML structure}]
<?xml version="1.0" encoding="UTF-8"?>
<pnml>
  <net id="net1" type="http://www.pnml.org/version-2009/grammar/ptnet">
    <page id="page1">
      <place id="p1">
        <name><text>Place 1</text></name>
        <initialMarking><text>1</text></initialMarking>
        <graphics><position x="100" y="100"/></graphics>
      </place>
      <transition id="t1">
        <name><text>Transition 1</text></name>
        <graphics><position x="200" y="100"/></graphics>
      </transition>
      <arc id="a1" source="p1" target="t1">
        <inscription><text>1</text></inscription>
      </arc>
    </page>
  </net>
</pnml>
\end{lstlisting}

\subsection{Graphical Information}

PNML files may optionally contain graphical information (positions, dimensions) in \texttt{<graphics>} elements. When this information is missing or invalid, the frontend applies the Sugiyama algorithm to generate a readable layout automatically.

\section{Web Technologies}
\label{sec:web-technologies}

The frontend is built using modern web technologies that enable rich, interactive single-page applications.

\subsection{Angular Framework}

Angular is a TypeScript-based web application framework developed by Google. Key concepts used in this project include:

\begin{itemize}
    \item \textbf{Components}: Reusable UI building blocks with encapsulated logic, templates, and styles
    \item \textbf{Services}: Singleton classes for shared state and business logic, injected via dependency injection
    \item \textbf{Observables (RxJS)}: Reactive programming patterns for handling asynchronous events and data streams
    \item \textbf{Modules}: Organizational units that group related components, services, and other code
\end{itemize}

\subsection{SVG for Petri Net Rendering}

Scalable Vector Graphics (SVG) is used for rendering Petri nets because:
\begin{itemize}
    \item Resolution-independent rendering at any zoom level
    \item DOM integration allows event handling on individual elements
    \item CSS styling and animations can be applied
    \item Native browser support without plugins
\end{itemize}

\subsection{Angular Material}

Angular Material provides pre-built UI components following Google's Material Design guidelines. The project uses Material components for:
\begin{itemize}
    \item Tab navigation (\texttt{mat-tab-group})
    \item Buttons and form controls
    \item Dialogs for warnings and information
    \item Icons and tooltips
\end{itemize}

\section{Related Concepts}
\label{sec:related-concepts}

% TODO: Add sections as needed for:
% - Docker containerization basics
% - Flask/REST API fundamentals
% - Any other concepts that need explanation




\chapter{<Weiteres Kapitel>}
\label{chap:weiteres-kapitel}


Natürlich sollten Sie den Überschriften sinnvoll wählen und nicht "`<Weiteres Kapitel>"' schreiben. Auch sollten Sie alle "`<"' und "`>"' in den Überschriften entfernen.

Dieses Kapitel ist das Haupt-Kapitel Ihrer Arbeit. Typischerweise gibt es davon auch mehrere. Haben Sie z.B. einen Algorithmus, Verfahren o.ä. entwickelt, implementiert und getestet, macht es vielleicht Sinn, in  Kapitel 3 den Algorithmus konzeptuell und unabhängig von der konkreten Programmierung zu beschreiben; welche Lösungsalternativen haben Sie ggf. abgewogen? In Kapitel 4 könnten Sie dann implementierungsspezifische Details des Algorithmus beschrieben, z.B. welche Programmiermuster Sie eingesetzt haben. In Kapitel 5 könnten Sie dann die Evaluation des Algorithmus und seiner Implementierung beschreiben.

Ist Ihre Arbeit über eine empirische Studie, könnten Sie in Kapitel 3 die Planung der Studie beschreiben, warum Sie denken, dass diese Studie eine bestimmte Hypothese oder Fragestellung möglichst gut beantworten wird. In Kapitel 4 könnten Sie dann die Ergebnisse der Studie präsentieren und interpretieren. In Kapitel 5 können Sie über die methodische Herangehensweise reflektieren.

Details dazu besprechen Sie am besten mit Ihrer Betreuerin oder Ihrem Betreuer.

\input{contents/06-verwandte-arbeiten.tex}

\input{contents/07-zusammenfassung-ausblick.tex}

\appendix

\chapter{Appendix: API Reference}
\label{chap:appendix-a}

This appendix provides a reference for the backend API endpoints used by the frontend.

\section{Simulation Endpoints}
\label{sec:simulation-endpoints}

\subsection{Simple Simulation}

\begin{description}
    \item[Endpoint:] \texttt{POST /simulation-petri-nets/simple-sim}
    \item[Description:] Executes a stochastic simulation of a Petri net
    \item[Request Body:] PNML content as string or file upload
    \item[Response:] JSON containing firing sequence and optional detailed log
\end{description}

\begin{lstlisting}[language=json, caption={Example response from simple-sim}]
{
  "firing_seq": ["t1", "t3", "t2", "t1", "t4"],
  "detailed_log": [
    {"step": 0, "marking": {"p1": 1, "p2": 0}},
    {"step": 1, "marking": {"p1": 0, "p2": 1}}
  ]
}
\end{lstlisting}

\subsection{Example Models}

\begin{description}
    \item[Endpoint:] \texttt{GET /simulation-petri-nets/example-models}
    \item[Description:] Retrieves a list of available example PNML files
    \item[Response:] Array of model names
\end{description}

\begin{description}
    \item[Endpoint:] \texttt{GET /simulation-petri-nets/example-models/<name>}
    \item[Description:] Retrieves the PNML content of a specific example model
    \item[Response:] PNML XML as string
\end{description}

\section{Planning Endpoints}
\label{sec:planning-endpoints}

\subsection{Default Parameters}

\begin{description}
    \item[Endpoint:] \texttt{GET /dto/planning}
    \item[Description:] Retrieves default parameter values for offshore scheduling
    \item[Response:] JSON object with parameter definitions and defaults
\end{description}

\subsection{Schedule Calculation}

\begin{description}
    \item[Endpoint:] \texttt{POST /offshore-plan/calculate-schedule}
    \item[Description:] Calculates an optimal schedule based on provided parameters
    \item[Request Body:] JSON with scenario and workforce configuration
    \item[Response:] Schedule with operation timings
\end{description}

\section{Configuration Files}
\label{sec:config-files}

\subsection{Angular Environment}

\texttt{src/environments/environment.ts}:

\begin{lstlisting}[language=TypeScript, caption={Development environment configuration}]
export const environment = {
  production: false,
  apiUrl: 'http://localhost:9040'
};
\end{lstlisting}

\subsection{Docker Compose}

\texttt{docker-compose.yml} for local development with backend:

\begin{lstlisting}[language=yaml, caption={Docker Compose configuration}]
version: '3.8'
services:
  frontend:
    build: .
    ports:
      - "8080:80"
    depends_on:
      - backend
  backend:
    image: l3s-offshore-backend:latest
    ports:
      - "9040:9040"
\end{lstlisting}


\bibliographystyle{abbrv}
%\bibliographystyle{alpha}
\bibliography{meinearbeit}

\backmatter


\end{document}